% Evans 6.4: English printed p.323 / PDF p.340; Chinese printed p.229 / PDF p.247.
\Needspace{65mm}
\section{Maximum Principles}
\label{sec:6-4}
Here we use nondivergence notation
\[
 Lu=-a^{ij}u_{ij}+b^iu_i+cu.
\]
Assume $U$ bounded, $A$ uniformly elliptic, and coefficients bounded and
continuous unless a stronger hypothesis is stated.
Solutions are classical, $u\in C^2(U)\cap C(\overline U)$.
These are different assumptions from the measurable-coefficient weak theory.

\subsection{Weak maximum principles (Evans \S6.4.1, Theorems 1--2)}
\begin{theorem}\label{thm:6-weak-max}
If $c=0$ and $Lu\leq0$, then
$\max_{\overline U}u=\max_{\partial U}u$.
If $c\geq0$ and $Lu\leq0$, then
\[
 \max_{\overline U}u\leq\max_{\partial U}u^+.
\]
In particular, $Lu\leq Lv$ and $u\leq v$ on $\partial U$ imply $u\leq v$
in $U$ when $c\geq0$.
\end{theorem}
\Needspace{45mm}
\PDEAnnotation{先固定算子符号}{本章主部是负号 $-a^{ij}u_{ij}$。因此 $Lu\leq0$ 控制上界；在内部极大点，$Du=0$、$D^2u\leq0$，主部反而非负。换符号时不等式也要一起换。}
\begin{expandedproof}
First let $c=0$. A strict subsolution $Lu<0$ cannot have an interior
maximum: at such a point $Du=0$, $D^2u$ is negative semidefinite,
and $-A:D^2u\geq0$, a contradiction.
For a non-strict subsolution use $u_\varepsilon=u+\varepsilon e^{\tau x_1}$.
Since
\[
 L(e^{\tau x_1})=(-a^{11}\tau^2+b^1\tau)e^{\tau x_1}<0
\]
for a fixed sufficiently large $\tau$, $u_\varepsilon$ has its maximum
on the boundary. Let $\varepsilon\downarrow0$.

For $c\geq0$, on the open set $D=\{u>0\}$ the operator
$K=L-c$ satisfies $Ku=Lu-cu\leq0$.
Apply the zero-potential result on $D$; boundary values on
$\partial D\cap U$ are zero, while the remaining values come from
$\partial U$. The proof uses compact maxima and does not require a
smooth level-set boundary. If $D$ is empty the conclusion is immediate.
Apply this assertion to $u-v$ for comparison.
\end{expandedproof}
Applying the theorem to $-u$ gives the corresponding minimum principles.
The restriction on $c$ is substantive: on $(0,\pi)$,
$u=\sin x$ solves $-u''-u=0$ with zero boundary values but is positive inside.

\Needspace{45mm}
\PDEAnnotation{弱最大值原理中的“弱”}{这里“弱”表示结论只控制边界极值，并非指 Sobolev 弱解。对弱解也有相应原理，但必须说明测试函数和系数条件，不能混用两个含义。}
For comparison, a direct weak energy argument is available in divergence
form when $b\in C^1(\overline U)$ and
$c-\tfrac12\operatorname{div}b\geq0$.
For a weak subsolution with $u^+\in H^1_0(U)$, testing with $u^+$ gives
\[
 0\geq B[u,u^+]
 =\int A Du^+\cdot Du^+
  +\int(c-\tfrac12\operatorname{div}b)(u^+)^2
 \geq\theta\|Du^+\|_2^2.
\]
Thus $u\leq0$. The drift integral was integrated by parts after writing
$u^+Du^+=D((u^+)^2)/2$.
This sufficient energy condition is distinct from $c\geq0$ in the
classical nondivergence argument.

\Needspace{35mm}
\subsection{Hopf's boundary lemma (Evans \S6.4.2)}
\begin{lemma}\label{lem:6-hopf}
Suppose an interior ball $B_r(y)\subset U$ touches $\partial U$ at $x_0$,
$u(x)<u(x_0)=M$ for all $x\in U$, and $Lu\leq0$.
Assume either $c=0$, or $c\geq0$ and $M\geq0$.
If $u$ is $C^1$ at $x_0$ in the inward-ball neighborhood, then
\[
 \partial_\nu u(x_0)>0,\qquad \nu=(x_0-y)/r.
\]
At a smooth tangency this is the outward normal to $U$.
\end{lemma}
\begin{expandedproof}
On $r/2<|x-y|<r$, put
$h(x)=e^{-\alpha|x-y|^2}-e^{-\alpha r^2}$.
It is positive inside and zero on the outer sphere. A direct calculation,
with $s=|x-y|$, gives
\[
 Lh\leq e^{-\alpha s^2}
 \bigl(-4\theta\alpha^2s^2+
       2\alpha\operatorname{tr}A+2\alpha|b|s+c\bigr)\leq0
\]
for sufficiently large $\alpha$; for $c=0$ the same bound applies.
On the inner sphere $u-M$ has a strictly negative maximum, so choose
$\varepsilon>0$ with $u-M+\varepsilon h\leq0$ there.
It is also nonpositive on the outer sphere.
Moreover $L(u-M+\varepsilon h)\leq-cM\leq0$.
Comparison on the annulus gives $u-M\leq-\varepsilon h$.
Along the radius $x_0-t\nu$ divide by $t>0$ and let $t\downarrow0$:
\[
 \partial_\nu u(x_0)\geq-\varepsilon\partial_\nu h(x_0)
 =2\varepsilon\alpha r e^{-\alpha r^2}>0.
\]
\end{expandedproof}
\Needspace{45mm}
\PDEAnnotation{Hopf 引理提供严格斜率}{最大值原理只说内部不能超过边界。Hopf 引理进一步说：若内部严格更小，函数沿外法向到达边界极大值时必须有正斜率；内球条件保证可放入显式屏障。}

\Needspace{35mm}
\subsection{Strong maximum principle (Evans \S6.4.2, Theorems 3--4)}
\begin{theorem}\label{thm:6-strong-max}
Let $U$ be connected. If $Lu\leq0$ and $c=0$, attainment of a maximum
at an interior point forces $u$ to be constant.
For $c\geq0$ the same holds for a nonnegative interior maximum.
Dually, a nonpositive interior minimum of a supersolution forces constancy.
\end{theorem}
\begin{expandedproof}
Let $M=\max_Uu$, attained inside, and suppose both sets
$C=\{u=M\}$ and $V=\{u<M\}$ are nonempty.
A path inside $U$ from $V$ to $C$ shows that their common relative
boundary contains a point in $U$.
Choose $y\in V$ sufficiently close to this point so that the ball with
radius $\operatorname{dist}(y,C)$ is contained in $U$.
Its interior is in $V$ and its boundary touches $C$ at some $x_0\in U$.
Hopf's lemma on this ball gives a strictly positive radial derivative
at $x_0$. But $x_0$ is an interior maximum of a differentiable function,
so $Du(x_0)=0$. This contradiction proves the assertion.
The minimum statement follows by replacing $u$ by $-u$.
\end{expandedproof}
In particular, for $c\geq0$, a nonzero nonnegative supersolution
$Lu\geq0$ is strictly positive in a connected domain.

% Expanded special-case proof of Evans 6.4.3, Theorem 5.
\Needspace{35mm}
\subsection{Harnack's inequality (Evans \S6.4.3, Theorem 5)}
\begin{theorem}[Smooth principal-part case proved here]\label{thm:6-harnack}
Suppose $A$ is smooth and uniformly elliptic and
$-a^{ij}u_{ij}=0$, $u\geq0$ in $U$. For connected $V\Subset U$,
\[
 \sup_Vu\leq C\inf_Vu.
\]
The constant depends on the geometry of $V$ inside $U$, ellipticity and
local coefficient bounds, including the derivatives used below, but not on $u$.
\end{theorem}
\Needspace{45mm}
\PDEAnnotation{Harnack 是定量正性}{强最大值原理说非负非零解处处为正；Harnack 还控制内区任意两点的比值。它需要方程等式，不能随意用于只有单边不等式的函数。}
\begin{expandedproof}
\paragraph{1. Take a logarithm to estimate relative oscillation.}
First take $u>0$, and put $v=\log u$ and $w=a^{ij}v_iv_j$.
Then
\begin{equation}\label{eq:6-log}
 -a^{ij}v_{ij}=w,\qquad \theta|Dv|^2\leq w\leq\Lambda|Dv|^2.
\end{equation}
Interior smoothness justifies the following differentiations:
rewrite the smooth nondivergence equation in divergence form and
apply \S\ref{sec:6-3}. On a fixed compact region,
\[
 w_{kl}=2a^{ij}v_{ikl}v_j+2a^{ij}v_{ik}v_{jl}+R_{kl},
 \quad
 |R_{kl}|\leq\varepsilon|D^2v|^2+C_\varepsilon|Dv|^2.
\]
The remainder consists of $D^2A\,Dv\,Dv$ and $DA\,D^2v\,Dv$.
Differentiate \eqref{eq:6-log} once to obtain
$-a^{kl}v_{ikl}=w_i+(\partial_i a^{kl})v_{kl}$.
Also
$a^{ij}a^{kl}v_{ik}v_{jl}\geq\theta^2|D^2v|^2$;
this is seen by diagonalizing the positive matrix $A$ at the point.
Absorbing coefficient remainders gives, for
$\mathcal M=-a^{kl}\partial_{kl}-2a^{ij}v_j\partial_i$,
\[
 \mathcal Mw\leq-\gamma|D^2v|^2+C|Dv|^2
                \leq-\gamma_1w^2+Cw.
\]
The second inequality uses $w= -A:D^2v\leq C|D^2v|$.

\paragraph{2. Localize and retain the drift--cutoff error.}
Take $\zeta\in C^\infty_c(U)$ with $0\leq\zeta\leq1$ and $\zeta=1$
on a chosen ball $V$. Let $z=\zeta^4w$.
If its maximum is positive, it occurs where $\zeta>0$, and
$\zeta Dw=-4wD\zeta$. At this point $Dz=0$ and
$\mathcal Mz=-A:D^2z\geq0$.
Expanding the product yields
\[
 0\leq-\gamma_1\zeta^4w^2+
        C\zeta^2w+C\zeta^3w^{3/2}.
\]
The last term comes from
$-2ADv\cdot D(\zeta^4)\,w$ and uses $|Dv|\leq C\sqrt w$.
It must be included because this drift depends on $v$.
Set $Y=\zeta^2w$ at this maximum point. Then
\[
 \gamma_1Y^2\leq C(Y+Y^{3/2}).
\]
For $Y>0$, divide by $Y$ and solve the resulting quadratic inequality
in $\sqrt Y$ to get $Y\leq C$. Hence $z\leq C$ at its maximum,
and $w\leq C$, $|Dv|\leq C$ on $V$.

\Needspace{45mm}
\PDEAnnotation{对数变换后的漂移项}{新算子的漂移含有 $D\log u$，不能预先当成已有一致界的系数。截断计算保留三次项，再用四次项吸收，才能保证最终常数不依赖解。}
\paragraph{3. Integrate the logarithmic gradient.}
For $x,y$ in the ball, integrate along their segment:
$|v(x)-v(y)|\leq C|x-y|\leq C\operatorname{diam}V$.
Exponentiation bounds $u(x)/u(y)$.
For a general connected $V\Subset U$, cover it by finitely many balls
whose slightly larger concentric balls remain inside $U$.
A finite chain of overlapping balls connects their values; multiplying
the finitely many comparison constants proves the result.
Finally apply the estimate to $u+\delta$, which solves the same equation
because there are no lower-order terms, and let $\delta\downarrow0$.
\end{expandedproof}


\PDEBodyDiagram{6.4}{Local positivity comparisons propagate along a chain.}
{Choose overlapping balls whose slightly larger balls still lie inside $U$. Comparing at overlap points transfers a local Harnack bound from $x$ to $y$. The resulting constant depends on the finite chain and its distance from the boundary, which is why connected compact interior regions occur in the theorem.}
{ch06-harnack-chain}

\paragraph{Scope of the source theorem.}
Evans states Harnack for the full operator with lower-order terms
and indicates extensions to bounded measurable coefficients.
The proof supplied in the source, and reconstructed above, is restricted
to smooth $A$ and $b=c=0$. The broader theorem is a further regularity
result, cited there to Gilbarg--Trudinger; it is not proved by the
logarithmic computation above. In the general case $u+\delta$ need not
solve $Lu=0$, and differentiating measurable coefficients is unavailable.